% Lamina 11: sesgo sostenido de 8 semanas (escenario 4), umbral p99.
% Datos: results/linea_base/rev23b/por_semana.csv, nivel 0.01, columna recall (fraccion de
% estacion-dias atacados con alarma en cada semana). 240 episodios aislados de sus vecinas,
% 30 calendarios del banco 2024-jul 2026. Detectores congelados (no se entreno nada).
% Detector final = Ens-Vb-X0.2596; PARAAD-SA = PARAAD.
% El ensamble con referencia propia (Ens-V0-R) se quito el 2026-11-23 a pedido del tesista
% porque estorbaba; esta en el mismo CSV (cae a ~0 desde la semana 3 en el escalon).
% Izquierda: escalon (sesgo constante desde el dia 1). Derecha: deriva de 56 dias que crece
% poco a poco hasta A. Tono claro = A +20 ppb, tono oscuro = A +40 ppb.
% Bajo cada semana de la deriva: sesgo real medio de esa semana en ppb
% (results/linea_base/rev23b/cambio_por_semana.csv, redondeado a entero).
% Franja crema: semanas 1-4, lo que aguanta el detector final antes de que su mediana de
% 60 dias absorba el sesgo (recall medio 0.76 con +20 ppb).
\ref{leyendaSesgo}\par
\begin{tikzpicture}[
    claro/.style={color=#1!70!white, mark options={fill=#1!70!white, solid}},  % tono medio de la paleta (magMedia)
    sesgo/.style={font=\sffamily\scriptsize, text=tenue, anchor=east},
  ]
\begin{groupplot}[
    group style={group size=2 by 1, horizontal sep=1.3cm, yticklabels at=edge left},
    ejes tesis,
    scale only axis, width=5.2cm, height=2.35cm,
    xmin=0.5, xmax=8.5, xtick={1,...,8},
    ymin=0, ymax=1.2, ytick={0,0.5,1},  % espacio arriba para el rotulo de la franja
    every axis plot/.append style={mark size=1.6pt},
    title style={font=\sffamily\small\bfseries, text=noche, yshift=-2pt},
    xlabel style={font=\sffamily\small},
    clip=false,
  ]
  % ---------------- Escalon ----------------
  \nextgroupplot[title={Escalón: sesgo constante},
                 ylabel={Recall semanal}, ylabel style={font=\sffamily\small},
                 xlabel={Semana},
                 legend to name=leyendaSesgo, legend columns=4,
                 legend style={/tikz/every even column/.append style={column sep=0.35cm}}]
  \addplot[forget plot, fill=resalte, draw=none, mark=none] coordinates {(0.5,0) (0.5,1.2) (4.5,1.2) (4.5,0)} \closedcycle;
  \node[font=\sffamily\scriptsize, text=ensX, anchor=north] at (axis cs:2.5,1.19) {$\approx$\,4 semanas};
  % +40 ppb (oscuro, con leyenda)
  \addplot[det ensX] coordinates {(1,0.977) (2,0.956) (3,0.919) (4,0.866) (5,0.568) (6,0.131) (7,0.052) (8,0.062)};
  \addlegendentry{Detector final}
  \addplot[det paraad] coordinates {(1,0.930) (2,0.945) (3,0.937) (4,0.928) (5,0.934) (6,0.937) (7,0.947) (8,0.933)};
  \addlegendentry{PARAAD-SA}
  % Muestras de tono para la leyenda (no dibujan nada)
  \addlegendimage{color=tenue!70!white, line width=1.6pt, mark=none}
  \addlegendentry{+20 ppb}
  \addlegendimage{color=texto, line width=1.6pt, mark=none}
  \addlegendentry{+40 ppb}
  % +20 ppb (claro)
  \addplot[det ensX, claro=ensX, forget plot] coordinates {(1,0.855) (2,0.845) (3,0.741) (4,0.607) (5,0.294) (6,0.078) (7,0.027) (8,0.037)};
  \addplot[det paraad, claro=paraad, forget plot] coordinates {(1,0.549) (2,0.562) (3,0.546) (4,0.531) (5,0.568) (6,0.564) (7,0.575) (8,0.577)};

  % ---------------- Deriva ----------------
  \nextgroupplot[title={Deriva: crece poco a poco},
                 xticklabels={{1\\1\\2}, {2\\4\\7}, {3\\6\\11}, {4\\8\\16},
                              {5\\10\\20}, {6\\13\\24}, {7\\15\\28}, {8\\17\\31}},
                 xticklabel style={align=center, font=\scriptsize},
                 xlabel={Semana y, debajo, sesgo real en ppb}]
  % +40 ppb (oscuro)
  \addplot[det ensX] coordinates {(1,0.024) (2,0.127) (3,0.364) (4,0.625) (5,0.781) (6,0.822) (7,0.874) (8,0.856)};
  \addplot[det paraad] coordinates {(1,0.038) (2,0.074) (3,0.177) (4,0.404) (5,0.670) (6,0.797) (7,0.895) (8,0.918)};
  % +20 ppb (claro)
  \addplot[det ensX, claro=ensX] coordinates {(1,0.017) (2,0.046) (3,0.086) (4,0.152) (5,0.208) (6,0.286) (7,0.379) (8,0.370)};
  \addplot[det paraad, claro=paraad] coordinates {(1,0.037) (2,0.038) (3,0.060) (4,0.092) (5,0.139) (6,0.234) (7,0.365) (8,0.524)};
  % Nombre de los renglones del sesgo real, a la izquierda de la semana 1
  \node[sesgo] at ([yshift=-0.58cm]axis cs:0.5,0) {+20:};
  \node[sesgo] at ([yshift=-0.86cm]axis cs:0.5,0) {+40:};
\end{groupplot}
\end{tikzpicture}
